\documentclass[12pt]{article}
\usepackage[margin=2.5cm]{geometry}
\usepackage{amsmath, amssymb}
\usepackage{graphicx}

\newcommand{\EE}{\mathbb{E}}
\newcommand{\PP}{\mathbb{P}}
\newcommand{\var}{{\rm Var}}
\newcommand{\invq}{\mbox{\reflectbox{?}}}

\begin{document}

\begin{center}
\begin{tabular}{@{}l r@{}}
\textbf{Probabilidad y Estad\'istica} & Universidad de la Rep\'ublica\\
Curso 2026 & Facultad de Ciencias\\
\end{tabular}
\end{center}

\vspace{1.2em}

\begin{center}
{\Large\bfseries Pra\'ctico 5. Teoremas de l\'imite}
\end{center}

\vspace{1.2em}

\begin{enumerate}

\item Sea $X_1, X_2, \dots$ una sucesi\'on de variables aleatorias independientes con $\EE(X_i)=\mu$ y $\var(X_i)=\sigma_i^2$. Demuestre que si $n^{-2}\sum_{i=1}^{n}\sigma_i^2 \to 0$, entonces $\bar{X}\to\mu$ en probabilidad.

\item Sean las $X_i$ como en el Ejercicio~1, pero ahora con $\EE(X_i)=\mu_i$ y $n^{-1}\sum_{i=1}^{n}\mu_i \to \mu$. Demuestre que $\bar{X}\to\mu$ en probabilidad.

\item Suponga que el n\'umero de reclamos de seguros, $N$, presentados en un a\'no tiene distribuci\'on de Poisson con $\EE(N)=10{,}000$. Utilice la aproximaci\'on normal de la Poisson para aproximar $\PP(N>10{,}200)$.

\item Suponga que el n\'umero de accidentes de tr\'ansito, $N$, en un per\'iodo de tiempo dado tiene distribuci\'on de Poisson con $\EE(N)=100$. Utilice la aproximaci\'on normal de la Poisson para hallar $\delta$ tal que
\[\PP(100-\delta<N<100+\delta)\approx 0{,}9.\]

\item Utilizando funciones generadoras de momentos, demuestre que, cuando $n\to\infty$, $p\to 0$ y $np\to\lambda$, la distribuci\'on binomial con par\'ametros $n$ y $p$ tiende a la distribuci\'on de Poisson.

\item Utilizando funciones generadoras de momentos, demuestre que, cuando $\alpha\to\infty$, la distribuci\'on gamma con par\'ametros $\alpha$ y $\lambda$, debidamente estandarizada, tiende a la distribuci\'on normal est\'andar.

\item Demuestre que si $X_n\to c$ en probabilidad y si $g$ es una funci\'on continua, entonces $g(X_n)\to g(c)$ en probabilidad.

\item Compare la funci\'on de distribuci\'on de la Poisson con la aproximaci\'on normal para (a) $\lambda=10$, (b) $\lambda=20$ y (c) $\lambda=40$.

\item Compare la funci\'on de distribuci\'on de la binomial con la aproximaci\'on normal para (a) $n=20$ y $p=0{,}2$, y (b) $n=40$ y $p=0{,}5$.

\item Se lanza un dado de seis caras 100 veces. Utilizando la aproximaci\'on normal, calcule la probabilidad de que salga un seis entre 15 y 20 veces. Calcule la probabilidad de que la suma de los valores de las 100 tiradas sea menor que 300.

\item Un esc\'eptico presenta el siguiente argumento para mostrar que debe haber un defecto en el teorema del l\'imite central: ``Sabemos que la suma de variables aleatorias de Poisson independientes sigue una distribuci\'on de Poisson cuyo par\'ametro es la suma de los par\'ametros de los sumandos. En particular, si se suman $n$ variables aleatorias de Poisson independientes, cada una con par\'ametro $n^{-1}$, la suma tiene una distribuci\'on de Poisson con par\'ametro 1. El teorema del l\'imite central dice que, cuando $n$ tiende a infinito, la distribuci\'on de la suma tiende a una distribuci\'on normal; pero la Poisson con par\'ametro 1 no es normal.'' \invq Qu\'e opina de este argumento?

\end{enumerate}

\vspace{2em}

{\small Los ejercicios fueron extra\'idos del libro:}
\medskip

{\small
John A. Rice, \emph{Mathematical Statistics and Data Analysis}, 2.\textsuperscript{a} ed., Wiley, 2006 (Secci\'on 5.4, Cap\'itulo 5: Teoremas l\'imite).
}

\end{document}
